\chapter{Stochastyczne równania różniczkowe: konstrukcja rozwiązań}
\label{ch:sde-existence}
\paragraph{Odsyłacze międzytomowe.} Numeryczne zastosowania SDE łączą się z aproksymacją z rozdziału~\ref{ch:sde-properties} oraz symulacyjnymi narzędziami tomu IV, rozdziału~\ref{IV-ch:monte-carlo}.


\section{Dlaczego równanie stochastyczne nie jest zwykłym równaniem różniczkowym?}
Rozważmy równanie $\dd X_t=b(t,X_t)\,\dd t+\sigma(t,X_t)\,\dd W_t$.
Gdyby $W$ miał gładkie trajektorie, można by próbować zastąpić
ostatni człon przez $\sigma(t,X_t)\dot W_t\,\dd t$.
Wiemy jednak z tomu II, że trajektorie procesu Wienera
są niemal na pewno nieróżniczkowalne. Ponadto całka Itô
została zdefiniowana przez domknięcie w przestrzeni procesów
przewidywalnych, a nie jako klasyczna całka wzdłuż trajektorii.
Zapis różniczkowy jest więc skrótem równania całkowego.

W teorii RÓR zadane pole wektorowe i punkt początkowy wyznaczają
trajektorię. Tu do tych danych dochodzi \emph{realizacja szumu}.
Pytamy osobno, czy rozwiązanie istnieje, czy jest jednoznaczne
przy ustalonym szumie, czy jego rozkład jest jednoznaczny i czy
może eksplodować w skończonym czasie. Rozróżnienia te nie są
formalnością: brak odpowiedniej regularności współczynników
może zniszczyć jednoznaczność nawet w równaniu deterministycznym.

\begin{center}
\begin{tikzpicture}[>=Latex,node distance=6mm]
\node[draw,rounded corners,align=center,fill=blue!5,minimum width=31mm] (data)
  {$X_0$, $b$, $\sigma$\\oraz proces $W$};
\node[draw,rounded corners,align=center,fill=orange!9,right=13mm of data,minimum width=37mm] (eq)
  {$X_t=X_0+\int_0^t b(s,X_s)\,\dd s$\\$\displaystyle+\int_0^t\sigma(s,X_s)\,\dd W_s$};
\node[draw,rounded corners,align=center,fill=green!8,right=13mm of eq,minimum width=32mm] (sol)
  {ciągły proces\\adaptowany $X$};
\draw[->,thick] (data)--(eq);
\draw[->,thick] (eq)--(sol);
\end{tikzpicture}
\par\smallskip\small Równanie stochastyczne jest warunkiem na \emph{cały proces}, nie wzorem na pochodną jego trajektorii.
\end{center}

\section{Definicje i dokładne znaczenie całek}
\begin{definicja}[Rozwiązanie silne na ustalonej przestrzeni]
\label{def:sde-strong}
Niech $T>0$ i niech $(\Omega,\calF,\Prob)$ będzie
przestrzenią probabilistyczną wyposażoną w filtrację
$(\calF_t)_{0\le t\le T}$ spełniającą warunki zwykłe.
Niech $W=(W^1,\ldots,W^m)$ będzie procesem Wienera
$m$-wymiarowym względem tej filtracji. Niech $X_0$ będzie
$\calF_0$-mierzalną zmienną o wartościach w $\R^d$.
Niech $b:[0,T]\times\R^d\to\R^d$ i
$\sigma:[0,T]\times\R^d\to\R^{d\times m}$
będą funkcjami borelowskimi. \emph{Silnym rozwiązaniem} równania
\[
\dd X_t=b(t,X_t)\,\dd t+\sigma(t,X_t)\,\dd W_t,
\qquad X_0\text{ zadane},
\tag{5.1}\label{eq:sde}
\]
nazywamy ciągły, adaptowany proces $X$ o wartościach w $\R^d$,
dla którego prawie na pewno
$\int_0^T (\norm{b(s,X_s)}+\norm{\sigma(s,X_s)}_{\mathrm F}^2)\,\dd s<\infty$
oraz, z prawdopodobieństwem jeden jednocześnie dla wszystkich $t$,
\[
X_t=X_0+\int_0^tb(s,X_s)\,\dd s
       +\int_0^t\sigma(s,X_s)\,\dd W_s.
\]
Druga całka jest sumą $m$ całek Itô po kolumnach macierzy $\sigma$.
\end{definicja}

\begin{uwaga}[Mierzalność i rola filtracji]
Ponieważ $X$ jest ciągły i adaptowany, jest przewidywalny.
Borelowskość $b$ i $\sigma$ zapewnia, że
$(t,\omega)\mapsto b(t,X_t(\omega))$ i
$(t,\omega)\mapsto\sigma(t,X_t(\omega))$ są przewidywalne,
a więc całki w definicji mają sens. Ważne jest, że nie wystarcza
samo istnienie procesu Wienera jako rodziny zmiennych losowych:
jego przyszłe przyrosty muszą być niezależne od aktualnej filtracji.
\end{uwaga}

\begin{definicja}[Jednoznaczność ścieżkowa i jednoznaczność co do prawa]
Dla ustalonych danych $b,\sigma$ mówimy o
\emph{jednoznaczności ścieżkowej}, gdy dowolne dwa rozwiązania
silne $X,Y$ zbudowane na tej samej filtrowanej przestrzeni,
z tym samym procesem Wienera i tym samym warunkiem początkowym
($X_0=Y_0$ prawie na pewno), są nierozróżnialne:
$\Prob(X_t=Y_t\ \text{dla wszystkich }t\in[0,T])=1$.
\emph{Jednoznaczność co do prawa} oznacza, że dowolne dwa
rozwiązania słabe o tym samym rozkładzie początkowym mają
ten sam rozkład na $C([0,T];\R^d)$.
\end{definicja}

\begin{definicja}[Rozwiązanie słabe]
Niech $\mu$ będzie miarą probabilistyczną na $\R^d$.
\emph{Rozwiązaniem słabym} \eqref{eq:sde} z rozkładem
początkowym $\mu$ nazywamy pewną filtrowaną przestrzeń
probabilistyczną, pewien $m$-wymiarowy proces Wienera $W$
względem jej filtracji oraz ciągły proces adaptowany $X$
spełniający równanie całkowe i $\Law(X_0)=\mu$.
W tej definicji przestrzeń probabilistyczna i szum nie są
a priori zadane. Słowo „słaby” nie oznacza, że równanie
obowiązuje tylko po scałkowaniu z funkcją testową.
\end{definicja}

\section{Podstawowe oszacowanie dla całek stochastycznych}
\begin{lemat}[Oszacowanie supremum]
\label{lem:sde-sup}
Niech $T>0$, $A$ będzie procesem progresywnie mierzalnym
o wartościach w $\R^d$, a $B$ przewidywalnym o wartościach
w $\R^{d\times m}$. Załóżmy, że
$\E\int_0^T(\norm{A_s}^2+\norm{B_s}_{\mathrm F}^2)\,\dd s<\infty$.
Wówczas dla $0\le t\le T$
\[
\E\sup_{r\le t}\norm{\int_0^r A_s\,\dd s+
\int_0^r B_s\,\dd W_s}^2
\le 2t\int_0^t\E\norm{A_s}^2\,\dd s
+8\int_0^t\E\norm{B_s}_{\mathrm F}^2\,\dd s.
\tag{5.2}\label{eq:sde-basic-est}
\]
\end{lemat}
\begin{proof}
Z nierówności $(u+v)^2\le2u^2+2v^2$ dla norm wektorowych
wystarczy osobno oszacować obie całki. Dla pierwszej,
przez Cauchy'ego--Schwarza, dla każdego $r\le t$ mamy
$\norm{\int_0^r A_s\,\dd s}^2
\le r\int_0^r\norm{A_s}^2\,\dd s
\le t\int_0^t\norm{A_s}^2\,\dd s$.
Dla drugiej, stosując nierówność Dooba w $L^2$
do kwadratu normy martyngału (lub do jego współrzędnych),
uzyskujemy
$\E\sup_{r\le t}\norm{\int_0^rB_s\,\dd W_s}^2
\le4\E\norm{\int_0^tB_s\,\dd W_s}^2$.
Izometria Itô daje $\E\norm{\int_0^tB_s\,\dd W_s}^2
=\E\int_0^t\norm{B_s}_{\mathrm F}^2\,\dd s$.
Po połączeniu oszacowań otrzymujemy tezę.
\end{proof}

\begin{uwaga}
Oszacowanie kontroluje \emph{całą trajektorię} przy pomocy
norm współczynników, a nie tylko wartość w ustalonej chwili.
Właśnie dlatego metoda Picarda zbiega w przestrzeni
ciągłych procesów, co jest silniejsze niż zbieżność
punktowa dla każdego $t$.
\end{uwaga}

\section{Istnienie i jednoznaczność przy globalnym warunku Lipschitza}
\begin{twierdzenie}[Twierdzenie o istnieniu i jednoznaczności silnego rozwiązania]
\label{th:sde-global}
Niech $T>0$, niech $W$ i filtracja będą jak w definicji
\ref{def:sde-strong}, a $X_0\in L^2(\Omega;\R^d)$ będzie
$\calF_0$-mierzalny. Załóżmy, że borelowskie współczynniki
$b:[0,T]\times\R^d\to\R^d$ i
$\sigma:[0,T]\times\R^d\to\R^{d\times m}$ spełniają:
\begin{enumerate}[label=(\roman*)]
\item dla pewnego $L\ge0$ i wszystkich $t\in[0,T]$ oraz
$x,y\in\R^d$ mamy
$\norm{b(t,x)-b(t,y)}+\norm{\sigma(t,x)-\sigma(t,y)}_{\mathrm F}
\le L\norm{x-y}$;
\item dla pewnego $K<\infty$ i wszystkich $t\in[0,T]$
mamy $\norm{b(t,0)}+\norm{\sigma(t,0)}_{\mathrm F}\le K$.
\end{enumerate}
Wówczas równanie \eqref{eq:sde} ma dokładnie jedno silne
rozwiązanie w sensie nierozróżnialności. Ponadto dla pewnej
stałej $C=C(T,L,K)$ zachodzi
\[
\E\sup_{0\le t\le T}\norm{X_t}^2
\le C(1+\E\norm{X_0}^2).
\tag{5.3}\label{eq:sde-moment}
\]
\end{twierdzenie}
\begin{proof}
\emph{Krok 1: liniowy wzrost.} Z założeń i nierówności
trójkąta otrzymujemy, dla pewnego $K_1=K_1(L,K)$,
$\norm{b(t,x)}^2+\norm{\sigma(t,x)}_{\mathrm F}^2
\le K_1(1+\norm{x}^2)$, jednostajnie względem $t$.

\emph{Krok 2: iteracje Picarda.} Ustalmy $X^{(0)}_t=X_0$
i dla $n\ge0$ zdefiniujmy
\[
X_t^{(n+1)}=X_0+\int_0^tb(s,X_s^{(n)})\,\dd s
+\int_0^t\sigma(s,X_s^{(n)})\,\dd W_s.
\tag{5.4}\label{eq:sde-picard}
\]
Przez indukcję wszystkie iteracje są ciągłymi procesami
adaptowanymi i należą do przestrzeni
$L^2(\Omega;C([0,T];\R^d))$; w kroku indukcyjnym
stosujemy liniowy wzrost i lemat \ref{lem:sde-sup}.
Nie postulujemy zatem istnienia całek przed wykazaniem
ich całkowalności.

\emph{Krok 3: szybkość zbieżności różnic.}
Niech $D_n(t)=\E\sup_{r\le t}
\norm{X_r^{(n+1)}-X_r^{(n)}}^2$.
Z lematu \ref{lem:sde-sup} i warunku Lipschitza,
po zastosowaniu nierówności $r\le T$, otrzymujemy
$D_n(t)\le C_0\int_0^t D_{n-1}(s)\,\dd s$
dla $n\ge1$ i $C_0=L^2(2T+8)$.
Pierwszy przyrost jest całkowalny:
$D_0(T)<\infty$ na mocy oszacowania zastosowanego
do $b(s,X_0),\sigma(s,X_0)$. Indukcja daje
\[
D_n(T)\le D_0(T)\frac{(C_0T)^n}{n!}.
\tag{5.5}\label{eq:picard-factorial}
\]
Zatem
$\sum_{n=0}^\infty\|X^{(n+1)}-X^{(n)}\|_{L^2(\Omega;C)}
\le\sqrt{D_0(T)}\sum_{n=0}^\infty(C_0T)^{n/2}/\sqrt{n!}<\infty$.
Przestrzeń $L^2(\Omega;C([0,T];\R^d))$ jest zupełna,
więc istnieje jej element $X$ będący granicą iteracji.
Możemy wybrać podciąg zbieżny jednostajnie niemal
na pewno; wynika to np. z nierówności Markowa i
lematu Borela--Cantellego. W tej wersji $X$ ma
ciągłe trajektorie i jest adaptowany.

\emph{Krok 4: przejście do granicy w równaniu.}
Z warunku Lipschitza i lematu \ref{lem:sde-sup}
wynika, że różnica między procesami całkowymi po
podstawieniu $X^{(n)}$ i $X$ zbiega do zera
w normie $L^2(\Omega;C)$, ponieważ jest ograniczona
przez $C_0\int_0^T\E\norm{X_s^{(n)}-X_s}^2\,\dd s$.
Granica spełnia więc równanie całkowe. Liniowy
wzrost oraz $X\in L^2(\Omega;C)$ gwarantują
wszystkie warunki całkowalności z definicji.

\emph{Krok 5: jednoznaczność.} Jeśli $X,Y$ są
rozwiązaniami o tym samym warunku początkowym i
przyrostach Wienera, oszacowanie z kroku 3 daje
$u(t):=\E\sup_{r\le t}\norm{X_r-Y_r}^2
\le C_0\int_0^t u(s)\,\dd s$.
Lemat Gronwalla implikuje $u(t)=0$ dla każdego $t$,
co oznacza nierozróżnialność.
Jeśli $X,Y$ są jedynie rozwiązaniami w sensie definicji
\ref{def:sde-strong}, bez założenia kwadratowej całkowalności
supremów, nie można stosować powyższego oszacowania wprost.
Dla $R>0$ wprowadzamy wspólny czas zatrzymania
$\tau_R=\inf\{t\in[0,T]:\norm{X_t}\vee\norm{Y_t}\ge R\}\wedge T$.
Ponieważ procesy są ciągłe, a $X_0=Y_0$ prawie na pewno,
na $[0,\tau_R]$ integrandy różnic są ograniczone przez
stałą zależną od $R$. Stosujemy lemat \ref{lem:sde-sup}
do różnicy całek zatrzymanych, uzyskując
$u_R(t):=\E\sup_{r\le t}\norm{X_{r\wedge\tau_R}-Y_{r\wedge\tau_R}}^2
\le C_0\int_0^t u_R(s)\,\dd s$.
Lemat Gronwalla daje $u_R(t)=0$. Dla prawie każdej trajektorii
ciągłe procesy $X,Y$ są ograniczone na $[0,T]$, więc
$\tau_R=T$ dla wszystkich dostatecznie dużych $R$.
Przechodząc do granicy po rosnącym ciągu promieni,
otrzymujemy nierozróżnialność na całym $[0,T]$.

\emph{Krok 6: ograniczenie momentów.}
Z równania i nierówności
$\norm{u+v+w}^2\le3(\norm u^2+\norm v^2+\norm w^2)$,
po zastosowaniu lematu \ref{lem:sde-sup} do obu części,
otrzymujemy
\[
\E\sup_{r\le t}\norm{X_r}^2
\le3\E\norm{X_0}^2+C_1\int_0^t(1+
\E\sup_{u\le s}\norm{X_u}^2)\,\dd s,
\]
gdzie $C_1$ zależy tylko od $T,L,K$.
Z lematu Gronwalla wynika \eqref{eq:sde-moment}.
\end{proof}

\begin{uwaga}[Dlaczego w dowodzie pojawia się silnia?]
Dla odwzorowania Picarda na całym przedziale $[0,T]$
nie musimy otrzymać kontrakcji w jednym kroku.
Każda kolejna iteracja wprowadza jednak dodatkowe
całkowanie po czasie. Objętość $n$-wymiarowego
simpleksu $0<s_n<\dots<s_1<T$ wynosi $T^n/n!$;
to właśnie źródło silni w \eqref{eq:picard-factorial}.
\end{uwaga}

\begin{zadanie}[Dlaczego regularność współczynników ma znaczenie?]
Rozważyć równanie $\dd X_t=2\sqrt{|X_t|}\,\dd t$
z warunkiem $X_0=0$, bez składnika szumowego.

\emph{Rozwiązanie.} Dla każdego $c\ge0$ proces
deterministyczny $X_t=(t-c)_+^2$ spełnia równanie:
dla $t\le c$ obie strony są zerowe, a dla $t>c$
mamy $\dot X_t=2(t-c)=2\sqrt{X_t}$. Istnieje
więc nieskończenie wiele rozwiązań. Współczynnik
$b(x)=2\sqrt{|x|}$ nie jest lokalnie Lipschitzowski
w zerze. Nawet bez szumu warunku tego nie można
w ogólności po prostu pominąć.
\end{zadanie}

\begin{lemat}[Mierzalna konstrukcja iteracji Picarda]
\label{lem:picard-measurable}
Niech współczynniki $b,\sigma$ spełniają założenia
twierdzenia \ref{th:sde-global}, a $\mu$ będzie rozkładem
na $\R^d$ o skończonym drugim momencie.
Dla ustalonego $T>0$ iteracje Picarda można skonstruować
tak, że prawo procesu $X^{(n)}$ jest jednoznacznie
wyznaczone przez prawo pary $(X_0,W)$ na
$\R^d\times C_0([0,T];\R^m)$, niezależnie od
filtrowanej przestrzeni probabilistycznej.
\end{lemat}
\begin{proof}
Dla $n=0$ proces jest stały $X_0$, więc wniosek jest oczywisty.
Załóżmy, że proces $X^{(n)}$ otrzymaliśmy jako mierzalny
funkcjonał pary $(X_0,W)$ (z dokładnością do zbioru miary
zero wyznaczonego przez jej prawo). Proces
$H_t=\sigma(t,X_t^{(n)})$ jest przewidywalny i należy
do $L^2(\dd t\otimes\Prob)$; analogicznie
$b(t,X_t^{(n)})$ jest całkowalny po czasie.
Dla $h=T/N$ i $k\ge1$ określamy
$H_k^{(h)}=h^{-1}\int_{(k-1)h}^{kh}H_s\,\dd s$,
a następnie proces schodkowy równy $H_k^{(h)}$
na $[kh,(k+1)h)$, z wartością zero na $[0,h)$.
Każda wartość $H_k^{(h)}$ jest $\calF_{kh}$-mierzalna;
więc przybliżenie jest przewidywalne i jego całka
Itô jest skończoną sumą
$\sum_k H_k^{(h)}(W_{t\wedge t_{k+1}}-W_{t\wedge t_k})$.
Jest to funkcjonał mierzalny pary $(X_0,W)$.

Dla dowolnego procesu z klasy $L^2(\dd t\otimes\Prob)$
przesunięte średnie na poprzednim przedziale zbiegają
do oryginalnego procesu w tej normie. Najpierw sprawdzamy
to dla funkcji schodkowych w czasie, a następnie dla
wszystkich funkcji $L^2$ przez ich gęstość i fakt, że
uśrednianie jest kontrakcją w $L^2$ (Jensen oraz Fubini).
Izometria Itô i nierówność maksymalna Dooba dają stąd
zbieżność całek w $L^2(\Omega;C([0,T]))$.
Dla całki względem czasu mierzalność funkcjonału
jest bezpośrednia z aproksymacji całek Lebesgue'a.
Wybór podciągu zbieżnego niemal na pewno pozwala
zidentyfikować całe $X^{(n+1)}$ jako mierzalny
funkcjonał tej samej pary. Wszystkie konstrukcje
odbywają się względem ustalonego wspólnego prawa
$(X_0,W)$, więc prawo iteracji nie zależy od
konkretnej realizacji przestrzeni probabilistycznej.
Indukcja kończy dowód.
\end{proof}

\begin{wniosek}[Jednoznaczność co do prawa]
\label{cor:sde-law}
Przy założeniach twierdzenia \ref{th:sde-global},
niech $\mu$ będzie rozkładem probabilistycznym
na $\R^d$ o skończonym drugim momencie.
Wówczas wszystkie rozwiązania słabe równania
\eqref{eq:sde} z $\Law(X_0)=\mu$ mają ten sam
rozkład na $C([0,T];\R^d)$.
\end{wniosek}
\begin{proof}
Dla każdego rozwiązania słabego $X_0$ jest
$\calF_0$-mierzalny, a $W$ jest procesem Wienera
względem tej filtracji. Cała trajektoria $W$
jest więc niezależna od $X_0$, a wspólne prawo pary
$(X_0,W)$ jest zawsze równe
$\mu\otimes\mu_W$, gdzie $\mu_W$ jest miarą Wienera.
Lemat \ref{lem:picard-measurable} daje mierzalną
konstrukcję kolejnych iteracji Picarda na podstawie
tej pary, a zatem ich prawo nie zależy od wyboru
rozwiązania słabego. W szczególności nie zakładamy
nieudowodnionej zależności całki Itô od całej trajektorii
szumu: uzyskujemy ją jako granicę jawnych sum
z przewidywalnymi współczynnikami. Na mocy
\eqref{eq:picard-factorial} iteracje zbiegają w
prawdopodobieństwie w przestrzeni trajektorii
do rozwiązania. Ponieważ rozkład granicy w
przestrzeni metrycznej jest wyznaczony przez rozkłady
przybliżeń, otrzymujemy tezę.
\end{proof}

\begin{przyklad}[Ten sam rozkład nie oznacza tych samych trajektorii]
Dla równania $\dd X_t=\dd W_t$, $X_0=0$,
jedynym silnym rozwiązaniem przy zadanym szumie jest
$X_t=W_t$. Na innych przestrzeniach możemy wziąć
różne ruchy Browna; odpowiednie rozwiązania nie są
funkcjami tego samego wyniku $\omega$, lecz mają
identyczne prawa na przestrzeni trajektorii.
To ilustruje różnicę między jednoznacznością ścieżkową
a jednoznacznością co do prawa.
\end{przyklad}

\section{Stabilność względem warunku początkowego}
\begin{twierdzenie}[Stabilność w normie średniokwadratowej]
\label{th:sde-stability}
Niech $b,\sigma$ spełniają założenia twierdzenia
\ref{th:sde-global}, a $X,Y$ będą rozwiązaniami
na tej samej przestrzeni, napędzanymi tym samym $W$,
z warunkami początkowymi $X_0,Y_0\in L^2$.
Wówczas dla pewnego $C=C(T,L)$
\[
\E\sup_{t\le T}\norm{X_t-Y_t}^2
\le C\E\norm{X_0-Y_0}^2.
\tag{5.6}\label{eq:sde-stability}
\]
\end{twierdzenie}
\begin{proof}
Odejmujemy równania całkowe. Ponieważ
$X_t-Y_t=X_0-Y_0+\int_0^t(b(s,X_s)-b(s,Y_s))\,\dd s
+\int_0^t(\sigma(s,X_s)-\sigma(s,Y_s))\,\dd W_s$,
z lematu \ref{lem:sde-sup} i nierówności
$\norm{u+v}^2\le2\norm u^2+2\norm v^2$
otrzymujemy
$u(t)\le2\E\norm{X_0-Y_0}^2+C_2\int_0^t u(s)\,\dd s$,
gdzie $u(t)=\E\sup_{r\le t}\norm{X_r-Y_r}^2$.
Lemat Gronwalla daje \eqref{eq:sde-stability}.
\end{proof}

\begin{przyklad}[Co mierzy oszacowanie stabilności?]
Jeśli dla tego samego procesu Wienera rozwiązujemy
dwa równania o identycznych współczynnikach, ale
z warunkami początkowymi oddalonymi o $\varepsilon$,
to średni kwadrat \emph{największej w czasie} różnicy
między trajektoriami jest nie większy od $C\varepsilon^2$.
Wynik nie mówi, że \emph{dwie niezależnie zaszumione}
trajektorie są bliskie. Sprzężenie polegające na
użyciu tego samego $W$ jest częścią założenia.
\end{przyklad}

\section{Warunki lokalne i zagadnienie eksplozji}
\begin{definicja}[Czas eksplozji]
Dla procesu ciągłego $X$ określonego do losowej chwili
$\zeta$ i wartościowanego w $\R^d$ definiujemy
$\tau_R=\inf\{t<\zeta:\norm{X_t}\ge R\}$
(przyjmujemy $\inf\varnothing=\infty$).
Jeżeli $\tau_R\uparrow\zeta$, to $\zeta$ nazywamy
\emph{czasem życia} procesu. Proces jest
\emph{nieeksplozywny} na $[0,T]$, jeśli $\Prob(\zeta>T)=1$.
\end{definicja}

\begin{twierdzenie}[Lokalny Lipschitz i liniowy wzrost]
\label{th:sde-local}
Niech $T>0$, niech $X_0\in L^2(\Omega;\R^d)$,
a $b,\sigma$ będą borelowskie i dla każdego $R>0$
istnieje $L_R<\infty$, takie że dla wszystkich
$t\in[0,T]$, $\norm x,\norm y\le R$ zachodzi
$\norm{b(t,x)-b(t,y)}+\norm{\sigma(t,x)-\sigma(t,y)}_{\mathrm F}
\le L_R\norm{x-y}$.
Załóżmy także istnienie $K$ z własnością
$\norm{b(t,x)}^2+\norm{\sigma(t,x)}_{\mathrm F}^2
\le K(1+\norm{x}^2)$ dla wszystkich $t,x$.
Wówczas \eqref{eq:sde} ma jedno silne rozwiązanie
na $[0,T]$, jest ono nieeksplozywne i spełnia
oszacowanie momentów analogiczne do \eqref{eq:sde-moment}.
\end{twierdzenie}
\begin{proof}
\emph{Obcięcie.} Niech $\rho_R(x)$ będzie rzutem
na domkniętą kulę promienia $R$; jest on $1$-Lipschitzowski
i równy identyczności w tej kuli. Definiujemy
$b_R(t,x)=b(t,\rho_R(x))$,
$\sigma_R(t,x)=\sigma(t,\rho_R(x))$.
Nowe współczynniki są globalnie Lipschitzowskie
ze stałą $L_R$ i spełniają ten sam warunek liniowego
wzrostu z tą samą stałą $K$, ponieważ
$\norm{\rho_R(x)}\le\norm x$.
Z twierdzenia \ref{th:sde-global} otrzymujemy
globalne rozwiązania $X^{(R)}$.

\emph{Zgodność rozwiązań.} Dla $S>R$ rozwiązania
$X^{(R)}$ i $X^{(S)}$ mają identyczne współczynniki,
gdy oba przebywają w kuli promienia $R$.
Stosujemy oszacowanie różnicy z dowodu jednoznaczności
do procesów zatrzymanych przy pierwszym opuszczeniu tej kuli
przez którykolwiek z nich. Ponieważ różnica początkowa
wynosi zero, Gronwall daje zgodność do tej chwili;
a więc również zgodność czasów pierwszego wyjścia.
Można zatem skleić lokalne rozwiązania na przedziale
przed granicą tych czasów.

\emph{Wykluczenie eksplozji.} Oszacowanie momentów z
kroku 6 poprzedniego dowodu zależy od stałej liniowego
wzrostu $K$, nie od stałej Lipschitza $L_R$.
Zatem $\sup_R\E\sup_{t\le T}\norm{X_t^{(R)}}^2
\le C(T,K)(1+\E\norm{X_0}^2)$.
Dla czasu wyjścia $\tau_R$ z kuli o promieniu $R$
wynika z nierówności Markowa
\[
\Prob(\tau_R\le T)\le
\frac{C(T,K)(1+\E\norm{X_0}^2)}{R^2}.
\]
Prawa strona dąży do zera. Stąd z prawdopodobieństwem
jeden lokalne rozwiązanie można przedłużyć na cały
$[0,T]$ i otrzymujemy jednoznaczność przez zgodność
obciętych równań.
\end{proof}

\begin{zadanie}[Eksplozja bez warunku liniowego wzrostu]
Rozwiązać równanie $\dd X_t=X_t^2\,\dd t$,
$X_0=x_0>0$, i wyznaczyć czas jego eksplozji.

\emph{Rozwiązanie.} Równanie jest deterministyczne.
Rozdzielenie zmiennych prowadzi do
$X_t=x_0/(1-x_0t)$ dla $t<1/x_0$.
W chwili $t\uparrow1/x_0$ rozwiązanie dąży do
nieskończoności, zatem $\zeta=1/x_0$. Funkcja
$b(x)=x^2$ jest lokalnie Lipschitzowska, ale nie
spełnia globalnego warunku liniowego wzrostu.
To pokazuje, dlaczego w poprzednim twierdzeniu
występują dwa odrębne założenia.
\end{zadanie}

\section{Rozwiązane modele: wykładniczy i liniowy}
\begin{zadanie}[Geometryczny ruch Browna: rozwiązanie i rozkład]
\label{zad:gbm}
Niech $a,c\in\R$ i $x_0>0$. Rozwiązać
\[
\dd X_t=aX_t\,\dd t+cX_t\,\dd W_t,
\qquad X_0=x_0,
\]
i wyznaczyć $\E X_t$, $\Var(X_t)$ oraz
$\Prob(X_t>K)$ dla ustalonego $K>0$ i $t>0$.

\emph{Rozwiązanie.} Ponieważ współczynniki są
Lipschitzowskie, rozwiązanie istnieje i jest jednoznaczne.
Szukamy postaci $X_t=x_0e^{Y_t}$, gdzie
$Y_t=(a-c^2/2)t+cW_t$. Wzór Itô dla funkcji
wykładniczej daje
$\dd X_t=X_t(a-c^2/2+c^2/2)\,\dd t+cX_t\,\dd W_t$.
A zatem
\[
X_t=x_0\exp\bigl((a-\tfrac12c^2)t+cW_t\bigr).
\tag{5.7}\label{eq:gbm-solution}
\]
Zmienne $W_t$ mają rozkład $\Normal(0,t)$, a więc
$X_t$ ma rozkład lognormalny (dla $c\ne0$).
Dzięki $\E e^{uW_t}=e^{u^2t/2}$ otrzymujemy
$\E X_t=x_0e^{at}$,
$\E X_t^2=x_0^2e^{(2a+c^2)t}$ oraz
$\Var(X_t)=x_0^2e^{2at}(e^{c^2t}-1)$.
Dla $c\ne0$, po standaryzacji zmiennej normalnej,
\[
\Prob(X_t>K)=1-\Phi\!\left(
\frac{\log(K/x_0)-(a-c^2/2)t}{|c|\sqrt t}\right),
\]
gdzie $\Phi$ jest dystrybuantą $\Normal(0,1)$.
Gdy $c=0$, proces jest deterministyczny i powyższe
prawdopodobieństwo wynosi $\ind{\{x_0e^{at}>K\}}$.
W szczególności $X_t>0$ prawie na pewno dla wszystkich $t$,
co nie jest oczywiste z samej postaci różniczkowej.
\end{zadanie}

\begin{zadanie}[Proces Ornsteina--Uhlenbecka]
\label{zad:ou}
Niech $\theta>0$, $m\in\R$, $\eta>0$ i
$X_0=x_0\in\R$. Rozwiązać
$\dd X_t=\theta(m-X_t)\,\dd t+\eta\,\dd W_t$
i obliczyć rozkład $X_t$ oraz $\Cov(X_s,X_t)$ dla $s\le t$.

\emph{Rozwiązanie.} Wprowadzamy czynnik całkujący
$e^{\theta t}$. Wzór na iloczyn Itô i zerowa
kowariacja funkcji czasu z $X$ dają
$\dd(e^{\theta t}(X_t-m))=\eta e^{\theta t}\,\dd W_t$.
Stąd
\[
X_t=m+(x_0-m)e^{-\theta t}
+\eta\int_0^t e^{-\theta(t-u)}\,\dd W_u.
\tag{5.8}\label{eq:ou-solution}
\]
Ostatnia całka jest gaussowska, ma średnią zero
i wariancję $\eta^2\int_0^t e^{-2\theta(t-u)}\,\dd u
=\eta^2(1-e^{-2\theta t})/(2\theta)$.
Dla $s\le t$ rozbijamy całkę w chwili $s$.
Przyrost od $s$ do $t$ jest niezależny od $\calF_s$,
a część wcześniejsza jest równa
$e^{-\theta(t-s)}(X_s-\E X_s)$. Dlatego
\[
\Cov(X_s,X_t)=\frac{\eta^2}{2\theta}
 e^{-\theta(t-s)}(1-e^{-2\theta s}).
\]
Z rozkładu widać, że gdy $t\to\infty$,
$X_t$ zbiega według rozkładu do
$\Normal(m,\eta^2/(2\theta))$.
Nie oznacza to, że $X_t$ zbiega prawie na pewno
do stałej $m$: w stanie równowagi szum nadal działa.
\end{zadanie}

\begin{przyklad}[Liniowe równanie o zmiennym dryfie]
Dla deterministycznych, ograniczonych funkcji $a,g$
i $h\in L^2([0,T])$ rozwiązanie
$\dd X_t=(a(t)X_t+g(t))\,\dd t+h(t)\,\dd W_t$
ma postać
\[
X_t=e^{A(t)}\left(X_0+\int_0^te^{-A(s)}g(s)\,\dd s
+\int_0^te^{-A(s)}h(s)\,\dd W_s\right),
\qquad A(t)=\int_0^t a(s)\,\dd s.
\]
Wystarczy zastosować wzór na iloczyn do
$e^{-A(t)}X_t$. Jeżeli $X_0$ jest deterministyczny,
proces jest gaussowski: jest sumą funkcji deterministycznej
i całki Wienera z deterministycznym integrandem.
\end{przyklad}

\section{Podsumowanie i granica teorii}
Wykazaliśmy istnienie i jednoznaczność przy globalnym
warunku Lipschitza, a następnie przy lokalnym warunku
Lipschitza i liniowym wzroście. Otrzymaliśmy także
stabilność i pierwsze modele jawne. Nie twierdzimy,
że współczynniki muszą być globalnie Lipschitzowskie,
aby SDE miało rozwiązanie: jest to użyteczne kryterium
wystarczające. W następnym rozdziale wykorzystamy
rozwiązania do budowania procesów Markowa, ich generatorów,
porównywania trajektorii i badania aproksymacji numerycznych.
% ============================================================================
% ROZDZIAL 6 — SDE: wlasnosci, generatory, porownanie, aproksymacja
% ============================================================================
